\ifdefined\BookMaster
\def\BookEnd{}
\else
\documentclass[a4paper,12pt,fleqn,openany,twoside,fontset=windows]{ctexbook}

\usepackage[fleqn]{amsmath}
\usepackage{amssymb,amsthm,mathrsfs}
\usepackage{geometry}
\usepackage[dvipsnames]{xcolor}
\usepackage{graphicx}
\usepackage{tikz}
\usetikzlibrary{calc}
\usepackage[most]{tcolorbox}
\usepackage{fancyhdr}
\usepackage{titlesec}
\usepackage{tocloft}
\usepackage{enumitem}
\usepackage{microtype}
\usepackage{pifont}
\usepackage{hyperref}
\usepackage{romannum}
\usepackage{mathtools}

\definecolor{bookblue}{HTML}{264653}
\definecolor{bookteal}{HTML}{4F7C82}
\definecolor{booklight}{HTML}{EDF4F3}
\definecolor{bookgray}{HTML}{5E686B}

\geometry{
    inner=2.7cm,
    outer=2.5cm,
    top=2.7cm,
    bottom=2.7cm,
    headheight=18pt,
    headsep=18pt,
    footskip=25pt
}
\setlength{\mathindent}{0pt}
\setlength{\parindent}{5pt}
\setlength{\parskip}{3pt plus 1pt minus 1pt}
\linespread{1.25}
\allowdisplaybreaks[3]
\AtBeginDocument{%
    \setlength{\parindent}{5pt}
    \setlength{\abovedisplayskip}{8pt plus 2pt minus 2pt}
    \setlength{\belowdisplayskip}{8pt plus 2pt minus 2pt}
    \setlength{\abovedisplayshortskip}{5pt plus 1pt minus 1pt}
    \setlength{\belowdisplayshortskip}{5pt plus 1pt minus 1pt}
}

\hypersetup{
    unicode=true,
    colorlinks=true,
    linkcolor=bookblue,
    citecolor=bookblue,
    urlcolor=bookteal,
    bookmarksopen=true,
    pdfauthor={Chen},
    pdftitle={数学分析习题整理}
}

\renewcommand{\chaptermark}[1]{%
    \edef\@tmp{第\thechapter 章\quad #1}%
    \expandafter\markboth\expandafter{\@tmp}{}%
}
\fancypagestyle{main}{%
    \fancyhf{}
    \fancyhead[LE]{\small\color{bookgray}\nouppercase{\leftmark}}
    \fancyhead[RO]{\small\color{bookgray}数学分析习题整理}
    \fancyfoot[C]{\color{bookblue}\thepage}
    \fancyfoot[R]{\small\bfseries Chen\;\; github.com/mathlover-1/math-analysis}
    \renewcommand{\headrulewidth}{0.4pt}
    \renewcommand{\footrulewidth}{0pt}
}
\fancypagestyle{plain}{%
    \fancyhf{}
    \fancyfoot[C]{\color{bookblue}\thepage}
    \fancyfoot[R]{\small\bfseries Chen\;\; github.com/mathlover-1/math-analysis}
    \renewcommand{\headrulewidth}{0pt}
    \renewcommand{\footrulewidth}{0pt}
}
\pagestyle{main}

\titleformat{\chapter}[display]
  {\normalfont\sffamily\bfseries\color{bookblue}}
  {\raggedleft\fontsize{46}{50}\selectfont\thechapter}
  {-18pt}
  {\titlerule[1.2pt]\vspace{10pt}\filright\Huge}
\titlespacing*{\chapter}{0pt}{-15pt}{36pt}

\renewcommand{\contentsname}{目\quad 录}
\renewcommand{\cfttoctitlefont}{\hfill\sffamily\bfseries\Huge\color{bookblue}}
\renewcommand{\cftaftertoctitle}{\hfill}
\setlength{\cftbeforetoctitleskip}{5pt}
\setlength{\cftaftertoctitleskip}{28pt}
\setlength{\cftbeforechapskip}{12pt}
\renewcommand{\cftchapfont}{\sffamily\bfseries\large\color{bookblue}}
\renewcommand{\cftchappagefont}{\sffamily\bfseries\large\color{bookblue}}
\renewcommand{\cftchapleader}{\color{bookteal}\cftdotfill{2.5}}

\newtcolorbox{ti}{
    enhanced,
    breakable,
    colback=booklight!55!white,
    colframe=bookteal!35!white,
    boxrule=0.45pt,
    boxsep=0pt,
    left=10pt,
    right=10pt,
    top=9pt,
    bottom=9pt,
    before skip=14pt,
    after skip=14pt,
    arc=2pt,
    outer arc=2pt,
    before upper={\parindent=0pt}
}

\newenvironment{booknotebody}{%
    \begin{center}
    \begin{minipage}{0.84\textwidth}
    \songti\zihao{-4}\linespread{1.45}\selectfont
    \setlength{\parindent}{2\ccwd}
    \setlength{\parskip}{0.75em}
}{%
    \end{minipage}
    \end{center}
}

\newenvironment{booknotelongbody}{%
    \begingroup
    \songti\zihao{-4}\linespread{1.45}\selectfont
    \setlength{\parindent}{2\ccwd}
    \setlength{\parskip}{0.75em}
    \begin{list}{}{%
        \setlength{\leftmargin}{0.08\textwidth}
        \setlength{\rightmargin}{0.08\textwidth}
        \setlength{\topsep}{0pt}
        \setlength{\partopsep}{0pt}
        \setlength{\parsep}{0pt}
        \setlength{\itemsep}{0pt}
    }
    \item\relax
}{%
    \end{list}
    \endgroup
}

\fancypagestyle{plainnowm}{%
    \fancyhf{}
    \fancyfoot[C]{\color{bookblue}\thepage}
    \renewcommand{\headrulewidth}{0pt}
    \renewcommand{\footrulewidth}{0pt}
}

\newcommand{\booknotetitle}[2][plain]{%
    \clearpage
    \phantomsection
    \addcontentsline{toc}{chapter}{#2}
    \markboth{#2}{}
    \thispagestyle{#1}
    \vspace*{1.2cm}
    \begin{center}
        {\sffamily\heiti\bfseries\fontsize{26}{34}\selectfont\color{bookblue}#2\par}
        \vspace{0.8em}
        {\color{bookteal}\rule{4.5cm}{0.8pt}\par}
    \end{center}
    \vspace{0.6cm}
}

\renewcommand{\proofname}{证明}
\renewcommand{\qedsymbol}{\color{bookteal}\ensuremath{\blacksquare}}

\title{数学分析习题整理}
\author{Chen}
\date{}

\makeatletter
% ==================== 旧版封面/封底设计备份 ====================
% 2026-08 重设计,旧代码用 \iffalse...\fi 封存,新版见下方。
\iffalse
\renewcommand{\maketitle}{%
    \begin{titlepage}
        \thispagestyle{empty}
        \setcounter{page}{0}
        \noindent\hspace*{-\parindent}\begin{tikzpicture}[remember picture,overlay]
            \fill[bookblue] (current page.north west) rectangle ([yshift=-4.2cm]current page.north east);
            \fill[booklight] ([yshift=-4.2cm]current page.north west) rectangle (current page.south east);
            \draw[bookteal!55,line width=0.7pt]
                ([xshift=2.4cm,yshift=4.3cm]current page.south west)
                .. controls +(2.5cm,1.8cm) and +(-2.8cm,-1.3cm) ..
                ([xshift=-2.4cm,yshift=7.2cm]current page.south east);
            \draw[bookteal!35,line width=0.7pt]
                ([xshift=2.5cm,yshift=6.2cm]current page.south west)
                .. controls +(3.0cm,-1.2cm) and +(-3.0cm,1.6cm) ..
                ([xshift=-2.3cm,yshift=3.8cm]current page.south east);
        \end{tikzpicture}
        \vspace*{5.6cm}
        \noindent\color{bookblue}\rule{\textwidth}{1.4pt}\par
        \vspace{8pt}
        {\centering\sffamily\bfseries\fontsize{32}{42}\selectfont\color{bookblue}\@title\par}
        \vspace{8pt}
        \noindent\color{bookteal}\rule{\textwidth}{0.6pt}\par
        \vspace{2.6cm}
        {\centering\Large\color{bookgray}\@author\par}
        \vfill
        {\centering\color{bookteal}\large
            $\displaystyle \lim_{n\to\infty} a_n$\qquad
            $\displaystyle \int_a^b f(x)\,\mathrm{d}x$\qquad
            $\displaystyle \nabla f(x,y)$\par
        }
        \vspace*{1.5cm}
    \end{titlepage}
}

\newcommand{\makebackcover}{%
    \begingroup
    \let\cleardoublepage\clearpage
    \begin{titlepage}
        \thispagestyle{empty}
        \noindent\hspace*{-\parindent}\begin{tikzpicture}[remember picture,overlay]
            \fill[bookblue] (current page.north west) rectangle (current page.south east);
            \fill[booklight,rounded corners=7pt]
                ([xshift=2.2cm,yshift=-3.0cm]current page.north west)
                rectangle
                ([xshift=-2.2cm,yshift=3.0cm]current page.south east);
            \fill[bookteal!80]
                (current page.north west) --
                ([xshift=5.3cm]current page.north west) --
                ([xshift=2.1cm,yshift=-4.4cm]current page.north west) --
                ([yshift=-6.2cm]current page.north west) -- cycle;
            \fill[bookteal!55]
                (current page.south east) --
                ([xshift=-5.8cm]current page.south east) --
                ([xshift=-2.4cm,yshift=4.2cm]current page.south east) --
                ([yshift=6.0cm]current page.south east) -- cycle;
            \foreach \radius in {0.75,1.15,1.55,1.95} {
                \draw[bookteal!55,line width=0.55pt]
                    ([xshift=-2.9cm,yshift=-4.4cm]current page.north east)
                    circle[radius=\radius cm];
            }
            \foreach \radius in {0.65,1.05,1.45} {
                \draw[booklight!65,line width=0.55pt]
                    ([xshift=2.4cm,yshift=3.6cm]current page.south west)
                    circle[radius=\radius cm];
            }
            \fill[bookteal!16,rounded corners=3pt]
                ([xshift=3.5cm,yshift=-7.1cm]current page.north west)
                rectangle
                ([xshift=-5.0cm,yshift=7.2cm]current page.south east);
            \fill[bookteal!30,rounded corners=3pt]
                ([xshift=4.4cm,yshift=-8.0cm]current page.north west)
                rectangle
                ([xshift=-4.1cm,yshift=8.1cm]current page.south east);
            \begin{scope}[shift={([xshift=-0.75cm]current page.center)}]
                \clip[rounded corners=3pt] (-5.25,-4.75) rectangle (5.25,4.75);
                \foreach \v in {-3.6,-3.0,...,3.6} {
                    \draw[bookteal!48,line width=0.45pt,opacity=0.72]
                        plot[domain=-3.6:3.6,samples=45,smooth,variable=\u]
                        ({\u+0.55*\v},{0.28*\u-0.35*\v+0.11*(\u*\u-\v*\v)});
                }
                \foreach \u in {-3.6,-3.0,...,3.6} {
                    \draw[bookblue!42,line width=0.42pt,opacity=0.62]
                        plot[domain=-3.6:3.6,samples=45,smooth,variable=\v]
                        ({\u+0.55*\v},{0.28*\u-0.35*\v+0.11*(\u*\u-\v*\v)});
                }
                \draw[bookteal!70,line width=0.75pt,opacity=0.82]
                    plot[domain=-3.6:3.6,samples=55,smooth,variable=\u]
                    ({\u},{0.28*\u+0.11*\u*\u});
                \draw[bookblue!60,line width=0.7pt,opacity=0.75]
                    plot[domain=-3.6:3.6,samples=55,smooth,variable=\v]
                    ({0.55*\v},{-0.35*\v-0.11*\v*\v});
            \end{scope}
            \draw[bookteal!45,line width=0.65pt]
                ([xshift=3.0cm,yshift=-8.2cm]current page.north west)
                .. controls +(2.2cm,2.0cm) and +(-2.0cm,-2.1cm) ..
                ([xshift=-3.0cm,yshift=9.0cm]current page.south east);
            \draw[bookteal!25,line width=0.55pt]
                ([xshift=3.1cm,yshift=8.0cm]current page.south west)
                .. controls +(2.7cm,-1.7cm) and +(-2.6cm,1.8cm) ..
                ([xshift=-3.1cm,yshift=-8.7cm]current page.north east);
            \node[anchor=west,text=bookteal!72,scale=0.78] at
                ([xshift=2.8cm,yshift=-4.0cm]current page.north west) {%
                $\displaystyle e^{\mathrm{i}\pi}+1=0$};
            \node[anchor=east,text=bookblue!62,scale=0.67] at
                ([xshift=-2.7cm,yshift=4.0cm]current page.south east) {%
                $\displaystyle
                f(x)=\sum_{n=0}^{\infty}\frac{f^{(n)}(a)}{n!}(x-a)^n$};
        \end{tikzpicture}
        \mbox{}
    \end{titlepage}
    \endgroup
}
\fi

% ==================== 新版封面(2026-08 重设计) ====================
% 设计要点:顶部深绿纵向渐变+波浪下缘+两层半透明过渡波(解决深浅绿生硬过渡);
% 深色区螺旋线/同心圆母题;标题+英文副题+作者居中;
% 中部 Fourier 部分和逼近方波曲线;底部浅色波纹(无文字)。
\renewcommand{\maketitle}{%
    \begin{titlepage}
        \thispagestyle{empty}
        \setcounter{page}{0}
        \noindent\hspace*{-\parindent}\begin{tikzpicture}[remember picture,overlay]
            % ---------- 底色 ----------
            \fill[booklight] (current page.north west) rectangle (current page.south east);

            % ---------- 顶部深色区:纵向渐变 + 波浪下缘 ----------
            \shade[top color=bookblue, bottom color=bookteal!55!bookblue]
                (current page.north west) -- (current page.north east) --
                ([yshift=-8.6cm]current page.north east)
                .. controls ([xshift=-5.2cm,yshift=-10.6cm]current page.north east)
                         and ([xshift=6.0cm,yshift=-8.0cm]current page.north west) ..
                ([yshift=-10.0cm]current page.north west) -- cycle;

            % ---------- 过渡层 1:半透明波浪 ----------
            \fill[bookteal!60!bookblue,opacity=0.42]
                ([yshift=-8.6cm]current page.north east)
                .. controls ([xshift=-5.2cm,yshift=-10.6cm]current page.north east)
                         and ([xshift=6.0cm,yshift=-8.0cm]current page.north west) ..
                ([yshift=-10.0cm]current page.north west) --
                ([yshift=-12.6cm]current page.north west)
                .. controls ([xshift=6.5cm,yshift=-11.0cm]current page.north west)
                         and ([xshift=-5.5cm,yshift=-13.2cm]current page.north east) ..
                ([yshift=-11.4cm]current page.north east) -- cycle;

            % ---------- 过渡层 2:更浅的波浪 ----------
            \fill[bookteal!35,opacity=0.35]
                ([yshift=-10.6cm]current page.north east)
                .. controls ([xshift=-6.0cm,yshift=-12.4cm]current page.north east)
                         and ([xshift=5.0cm,yshift=-9.8cm]current page.north west) ..
                ([yshift=-11.8cm]current page.north west) --
                ([yshift=-13.8cm]current page.north west)
                .. controls ([xshift=5.5cm,yshift=-12.6cm]current page.north west)
                         and ([xshift=-6.5cm,yshift=-14.6cm]current page.north east) ..
                ([yshift=-12.8cm]current page.north east) -- cycle;

            % ---------- 深色区装饰:右上螺旋线 ----------
            \begin{scope}[shift={([xshift=-3.4cm,yshift=-4.3cm]current page.north east)}]
                \draw[booklight!45!bookteal,opacity=0.55,line width=0.6pt]
                    plot[domain=0:1440,samples=220,smooth,variable=\t]
                    ({0.0028*\t*cos(\t)},{0.0028*\t*sin(\t)});
            \end{scope}

            % ---------- 深色区装饰:左下同心圆 ----------
            \foreach \r in {0.7,1.25,1.8,2.35,2.9}{
                \draw[booklight!35!bookteal,opacity=0.5,line width=0.5pt]
                    ([xshift=0.4cm,yshift=-7.9cm]current page.north west) circle[radius=\r cm];
            }

            % ---------- 主标题 ----------
            \node[anchor=center,text=booklight] at
                ([yshift=-4.5cm]current page.north)
                {\sffamily\bfseries\fontsize{40}{46}\selectfont \@title};

            % ---------- 分隔短线 ----------
            \draw[booklight!65,line width=1.2pt]
                ([xshift=-2.6cm,yshift=-5.75cm]current page.north) --
                ([xshift=2.6cm,yshift=-5.75cm]current page.north);

            % ---------- 英文副题 ----------
            \node[anchor=center,text=booklight!70!bookteal] at
                ([yshift=-6.55cm]current page.north)
                {\itshape\large Mathematical Analysis \textperiodcentered\
                 A Collection of Exercises};

            % ---------- 作者(紧随标题下方) ----------
            \node[anchor=center,text=booklight!88] at
                ([yshift=-8.2cm]current page.north)
                {\sffamily\large \@author\quad 整理};

            % ---------- 中部:Fourier 部分和逼近方波 ----------
            \begin{scope}[shift={([xshift=-5.2cm,yshift=-18.1cm]current page.north)}]
                % 坐标轴
                \draw[bookgray!55,line width=0.5pt,->] (-0.5,0) -- (10.9,0);
                \draw[bookgray!55,line width=0.5pt,->] (0,-1.75) -- (0,1.95);
                % 目标方波(虚线)
                \draw[bookgray!75,line width=0.6pt,dashed]
                    (0,0.92) -- (5,0.92) (5,-0.92) -- (10,-0.92);
                % N=1
                \draw[bookteal!45,line width=0.7pt]
                    plot[domain=0:360,samples=90,smooth,variable=\x]
                    ({\x/36},{1.18*sin(\x)});
                % N=3
                \draw[bookteal!80,line width=0.8pt]
                    plot[domain=0:360,samples=120,smooth,variable=\x]
                    ({\x/36},{1.18*(sin(\x)+sin(3*\x)/3)});
                % N=7
                \draw[bookblue,line width=0.95pt]
                    plot[domain=0:360,samples=150,smooth,variable=\x]
                    ({\x/36},{1.18*(sin(\x)+sin(3*\x)/3+sin(5*\x)/5+sin(7*\x)/7)});
                % 刻度标注
                \node[anchor=north east,text=bookgray!80,scale=0.62] at (-0.12,0) {$0$};
                \node[anchor=north,text=bookgray!80,scale=0.62] at (5,0) {$\pi$};
                \node[anchor=north,text=bookgray!80,scale=0.62] at (10,0) {$2\pi$};
            \end{scope}
            % 图注
            \node[anchor=center,text=bookgray,scale=0.8] at
                ([yshift=-20.5cm]current page.north)
                {$S_N(x)=\displaystyle\sum_{k=1}^{N}\frac{\sin\bigl((2k-1)x\bigr)}{2k-1}$
                 \ 逐步逼近方波};

            % ---------- 底部浅色波纹(无文字,与顶部波浪呼应) ----------
            \fill[bookteal!22,opacity=0.5]
                ([yshift=2.4cm]current page.south west)
                .. controls ([xshift=6.0cm,yshift=3.9cm]current page.south west)
                         and ([xshift=-6.5cm,yshift=1.0cm]current page.south east) ..
                ([yshift=2.9cm]current page.south east) --
                (current page.south east) -- (current page.south west) -- cycle;
            \fill[bookteal!40!bookblue,opacity=0.28]
                ([yshift=1.3cm]current page.south west)
                .. controls ([xshift=5.5cm,yshift=2.7cm]current page.south west)
                         and ([xshift=-5.5cm,yshift=0.3cm]current page.south east) ..
                ([yshift=1.9cm]current page.south east) --
                (current page.south east) -- (current page.south west) -- cycle;
        \end{tikzpicture}
        \mbox{}
    \end{titlepage}
}

% ==================== 新版封底(纯背景,2026-08 重设计) ====================
% 设计要点:对角渐变底+角部半透明叠层;中下部外摆线(玫瑰形)为视觉主体;
% 底部两层半透明白色波浪;无任何文字。
\newcommand{\makebackcover}{%
    \begingroup
    \let\cleardoublepage\clearpage
    \begin{titlepage}
        \thispagestyle{empty}
        \noindent\hspace*{-\parindent}\begin{tikzpicture}[remember picture,overlay]
            % ---------- 底色:对角渐变 ----------
            \shade[top color=bookblue, bottom color=bookteal!48!bookblue,
                   shading angle=-30]
                (current page.north west) rectangle (current page.south east);

            % ---------- 角部半透明叠层 ----------
            \fill[white,opacity=0.05]
                ([xshift=2.0cm,yshift=1.0cm]current page.north east) circle[radius=7.2cm];
            \fill[white,opacity=0.06]
                ([xshift=-0.5cm,yshift=-1.5cm]current page.north east) circle[radius=3.6cm];

            % ---------- 左下同心圆 ----------
            \foreach \r in {1.0,1.75,2.5,3.25,4.0,4.75}{
                \draw[booklight!40,opacity=0.55,line width=0.55pt]
                    ([xshift=-0.6cm,yshift=-0.4cm]current page.south west) circle[radius=\r cm];
            }
            % ---------- 右下螺旋线 ----------
            \begin{scope}[shift={([xshift=-2.9cm,yshift=3.2cm]current page.south east)}]
                \draw[booklight!45!bookteal,opacity=0.5,line width=0.55pt]
                    plot[domain=0:1260,samples=200,smooth,variable=\t]
                    ({0.0036*\t*cos(\t)},{0.0036*\t*sin(\t)});
            \end{scope}
            % ---------- 左上斜线组 ----------
            \foreach \i in {0,1,2}{
                \draw[bookteal!60!booklight,opacity=0.5,line width=0.6pt]
                    ([xshift={1.6cm+\i*0.45cm}]current page.north west) --
                    ([yshift={-1.6cm-\i*0.45cm}]current page.north west);
            }

            % ---------- 中下部视觉主体:外摆线(玫瑰形) ----------
            \begin{scope}[shift={([yshift=-16.6cm]current page.north)}]
                % 外圈
                \draw[booklight!35,opacity=0.6,line width=0.6pt]
                    (0,0) circle[radius=4.9cm];
                \draw[booklight!22,opacity=0.5,line width=0.5pt]
                    (0,0) circle[radius=5.25cm];
                % 主玫瑰线
                \draw[booklight!55!bookteal,opacity=0.8,line width=0.85pt]
                    plot[domain=0:1080,samples=400,smooth,variable=\t]
                    ({0.335*(8*cos(\t)-5*cos(2.6667*\t))},
                     {0.335*(8*sin(\t)-5*sin(2.6667*\t))});
                % 内层玫瑰线
                \draw[booklight!38!bookteal,opacity=0.6,line width=0.6pt]
                    plot[domain=0:1080,samples=400,smooth,variable=\t]
                    ({0.21*(8*cos(\t)-5*cos(2.6667*\t))},
                     {0.21*(8*sin(\t)-5*sin(2.6667*\t))});
            \end{scope}

            % ---------- 底部半透明波浪 ----------
            \fill[white,opacity=0.06]
                ([yshift=4.6cm]current page.south west)
                .. controls ([xshift=6.5cm,yshift=6.4cm]current page.south west)
                         and ([xshift=-6.0cm,yshift=2.6cm]current page.south east) ..
                ([yshift=5.2cm]current page.south east) --
                (current page.south east) -- (current page.south west) -- cycle;
            \fill[booklight,opacity=0.08]
                ([yshift=2.6cm]current page.south west)
                .. controls ([xshift=5.5cm,yshift=4.3cm]current page.south west)
                         and ([xshift=-6.5cm,yshift=1.0cm]current page.south east) ..
                ([yshift=3.3cm]current page.south east) --
                (current page.south east) -- (current page.south west) -- cycle;
        \end{tikzpicture}
        \mbox{}
    \end{titlepage}
    \endgroup
}
\makeatother

% 页脚右下方常态化水印(署名 + 仓库地址,加粗黑色),已在 main/plain 分页样式中定义,所有页面均生效.

\begin{document}
\mainmatter
\pagestyle{main}
\setcounter{chapter}{0}
\def\BookEnd{\end{document}}
\fi

\chapter{数列}

% 在此处使用 ti 环境添加题目,例如:
% \begin{ti}
% 证明:若数列 $\{a_n\}$ 单调有界,则必收敛.
% \end{ti}
%
% 支持行间公式:
% \begin{ti}
% 求极限
% \[
%   \lim_{n \to \infty} \frac{a_n}{b_n}
% \]
% \end{ti}

\begin{ti}
【题1】 设$a_n=\left[\dfrac{(2n-1)!!}{(2n)!!}\right]^2$,证明:$\{a_n\}$收敛于0.
\end{ti}

\textbf{证明}:在不使用斯特林公式的情况下,对任意$\displaystyle n\geq1$,有 \\

$\displaystyle\prod_{k=1}^{n}\dfrac{(2k-1)(2k+1)}{(2k)^2}=\dfrac{1\cdot3^2\cdot5^2\cdots(2n-1)^2(2n+1)}{2^2\cdot4^2\cdots(2n)^2}=(2n+1)a_n,$ \\

又对每个$\displaystyle k\geq1$,有$\displaystyle 0<\dfrac{(2k-1)(2k+1)}{(2k)^2}=1-\dfrac{1}{4k^2}<1,$ \\

故$\displaystyle 0<(2n+1)a_n<1$,从而$\displaystyle 0<a_n<\dfrac{1}{2n+1}\to0,$ \\

由夹逼准则,$\displaystyle\lim_{n\to\infty}a_n=0.$

\begin{ti}
【题2】 $a_1> 0,\quad a_{n+1}=a_n+\dfrac{1}{a_n}\quad \text{证明}:\displaystyle\lim_{n\to\infty}\dfrac{a_n}{\sqrt{2n}}=1$
\end{ti}

\textbf{证明}: 易得$a_n> 0$,且$a_{n+1}> a_n,$ \\

若$\{a_n\}$收敛,设收敛于$a,\;a\in\mathbb{R},\text{则}a=a+\dfrac{1}{a}$,无解, \\

若$\displaystyle\{a_n\}$有上界,则由单调收敛定理必收敛,与上面的结论矛盾. \\

故$\displaystyle\{a_n\}$无上界;又$\displaystyle\{a_n\}$单调递增,所以$\displaystyle a_n\to+\infty$,则$\displaystyle\lim_{n\to\infty}\dfrac{1}{a_n}=0,$ \\

$\because\displaystyle\lim_{n\to\infty}\dfrac{a_{n+1}^2-a_n^2}{2(n+1)-2n}=\lim_{n\to\infty}\dfrac{\frac{1}{a_n^2}+2}{2}=1,$ \\

且$b_n=2n$为严格单调递增且发散到$+\infty$的数列, \\

由\textit{Stolz}定理,$\quad\displaystyle\lim_{n\to\infty}\dfrac{a_n^2}{2n}=1,$ \\

$\therefore\displaystyle\lim_{n\to\infty}\dfrac{a_n}{\sqrt{2n}}=1.$ \\

\textbf{一点思考}:关于为什么要平方,我的思考是:若有$\displaystyle\lim_{n\to\infty}(b_{n+1}-b_n)=0,$ \\

$\text{则由Stolz公式},\displaystyle\lim_{n\to\infty}\dfrac{b_n}{n}=0$, 而这里$a_{n+1}-a_n\rightarrow 0,$ \\

$\text{但是题目要求}\dfrac{a_n}{\sqrt{2n}}\rightarrow 1$,这说明$a_{n+1}-a_n\rightarrow 0$的条件不够强, \\

应该研究$b_n=a_n^2,\text{研究}b_{n+1}-b_n$.

\begin{ti}
【题3】 设$0< x_i\leq \frac{1}{2}\quad i=1,2,...,n,\text{证明}:  $ \\

$\dfrac{\displaystyle\prod_{i=1}^n x_i}{(\displaystyle\sum_{i=1}^n x_i)^n}\leq
\dfrac{\displaystyle\prod_{i=1}^n (1-x_i)}{[\displaystyle\sum_{i=1}^n (1-x_i)]^n}$
\end{ti}

\textbf{证明}:$\;n=1\text{时,上式显然成立},$ \\

$n=2\text{时,即证}\dfrac{x_1x_2}{(x_1+x_2)^2}\leq\dfrac{(1-x_1)(1-x_2)}{(2-x_1-x_2)^2},$ \\

$\text{即证}\quad x_1x_2\left[4+(x_1+x_2)^2-4(x_1+x_2)\right]\leq(x_1+x_2)^2\left[1+x_1x_2-(x_1+x_2)\right],$ \\

$\text{即证}\quad 4x_1x_2\left[1-(x_1+x_2)\right]\leq (x_1+x_2)^2\left[1-(x_1+x_2)\right],$ \\

$\because 1-(x_1+x_2)\geq 0,\text{即证}\quad 4x_1x_2\leq (x_1+x_2)^2\quad\text{易得上式成立},$ \\

$\text{下面使用向前向后证明法},$ \\

$\text{向前:}\quad n\rightarrow 2n,$ \\

$\begin{cases}
    &\dfrac{\displaystyle\prod_{i=1}^n x_i}{(\displaystyle\sum_{i=1}^n x_i)^n}\leq
\dfrac{\displaystyle\prod_{i=1}^n (1-x_i)}{[\displaystyle\sum_{i=1}^n (1-x_i)]^n} \text{\ding{172}} \\
    &\dfrac{\displaystyle\prod_{i=n+1}^{2n} x_i}{(\displaystyle\sum_{i=n+1}^{2n} x_i)^n}\leq
\dfrac{\displaystyle\prod_{i=n+1}^{2n} (1-x_i)}{[\displaystyle\sum_{i=n+1}^{2n} (1-x_i)]^n} \text{\ding{173}} \\
\end{cases}$
\[\resizebox{0.70\textwidth}{!}{$\text{\ding{172}}\times\text{\ding{173}}:\dfrac{\displaystyle\prod_{i=1}^{2n} x_i}{(\displaystyle\sum_{i=1}^n x_i\cdot\sum_{i=n+1}^{2n} x_i)^n}\leq\dfrac{\displaystyle\prod_{i=1}^{2n} (1-x_i)}{[\displaystyle\sum_{i=1}^n (1-x_i)\cdot\sum_{i=n+1}^{2n} (1-x_i)]^n}$}\] \\
$\text{设}A=\displaystyle\sum_{i=1}^n x_i\quad B=\displaystyle\sum_{i=n+1}^{2n} x_i\quad\text{则}A\leq\frac{n}{2},B\leq\frac{n}{2},$ \\

$\therefore\dfrac{\displaystyle\prod_{i=1}^{2n} x_i}{\displaystyle\prod_{i=1}^{2n} (1-x_i)}\leq\dfrac{(AB)^n}{(n-A)^n(n-B)^n},$ $\text{令}A'=\dfrac{A}{n},B'=\dfrac{B}{n},$ \\

$\therefore\dfrac{(AB)^n}{(n-A)^n(n-B)^n}=\left[\dfrac{A'B'}{(1-A')(1-B')}\right]^n\quad \text{由}n=2\text{情形可知},$ \\

$\left[\dfrac{A'B'}{(1-A')(1-B')}\right]^n\leq\dfrac{(A'+B')^{2n}}{(2-A'-B')^{2n}}=\dfrac{(A+B)^{2n}}{(2n-A-B)^{2n}},$ \\

即$\dfrac{\displaystyle\prod_{i=1}^{2n} x_i}{\displaystyle\prod_{i=1}^{2n} (1-x_i)}\leq\dfrac{(\displaystyle\sum_{i=1}^{2n} x_i)^{2n}}{(2n-\displaystyle\sum_{i=1}^{2n} x_i)^{2n}}=\dfrac{(\displaystyle\sum_{i=1}^{2n} x_i)^{2n}}{[\displaystyle\sum_{i=1}^{2n} (1-x_i)]^{2n}},$ \\
 
向后:$n\rightarrow n-1$,  假设对$n$成立, $\because\dfrac{1}{n-1}\displaystyle\sum_{i=1}^{n-1} x_i\leq\dfrac{1}{2}\quad\text{由}x_n\text{的任意性},$ \\

不妨取$x_n=\dfrac{1}{n-1}\displaystyle\sum_{i=1}^{n-1} x_i,$ \\

设$S_n=\displaystyle\sum_{i=1}^{n} x_i,\quad\text{则}x_n=\dfrac{S_{n-1}}{n-1},\quad S_n=S_{n-1}+x_n=S_{n-1}\cdot\dfrac{n}{n-1},$ \\

$\therefore\dfrac{\displaystyle\prod_{i=1}^n x_i}{(S_n)^n}\leq\dfrac{\displaystyle\prod_{i=1}^{n} (1-x_i)}{(n-S_n)^n},$ \\

$\therefore\dfrac{\displaystyle\prod_{i=1}^{n-1} x_i\cdot\dfrac{S_{n-1}}{n-1}}{S_{n-1}^n\cdot\frac{n^n}{(n-1)^n}}\leq\dfrac{\displaystyle\prod_{i=1}^{n-1} (1-x_i)\cdot(1-\dfrac{S_{n-1}}{n-1})}{n^n(1-\dfrac{S_{n-1}}{n-1})^n},$ \\

$\therefore\dfrac{\displaystyle\prod_{i=1}^{n-1} x_i}{(\dfrac{S_{n-1}}{n-1})^{n-1}}\leq\dfrac{\displaystyle\prod_{i=1}^{n-1} (1-x_i)}{(1-\dfrac{S_{n-1}}{n-1})^{n-1}},$ \\

$\therefore\dfrac{\displaystyle\prod_{i=1}^{n-1} x_i}{(\displaystyle\sum_{i=1}^{n-1} x_i)^{n-1}}\leq\dfrac{\displaystyle\prod_{i=1}^{n-1} (1-x_i)}{\left[\displaystyle\sum_{i=1}^{n-1} (1-x_i)\right]^{n-1}},$ \\

又$n=2$时命题成立, \\
由向前-向后归纳法可知 $\dfrac{\displaystyle\prod_{i=1}^{n} x_i}{(\displaystyle\sum_{i=1}^{n} x_i)^n}\leq\dfrac{\displaystyle\prod_{i=1}^{n} (1-x_i)}{\left[\displaystyle\sum_{i=1}^{n} (1-x_i)\right]^n}.$

\begin{ti}
【题4】设$a_n=\sqrt{1+\sqrt{2+\sqrt{...+\sqrt{n}}}}\quad $证明:$\{a_n\}$收敛
\end{ti}

\textbf{证明}: 易得$\{a_n\}$单调递增, \\

对固定的$\displaystyle n$,定义$\displaystyle x_{n,n}=\sqrt n,\quad x_{k,n}=\sqrt{k+x_{k+1,n}}\quad(k=n-1,\ldots,1)$,则$\displaystyle a_n=x_{1,n},$ \\

下面用反向归纳法证明$\displaystyle x_{k,n}\leq\sqrt{k}+1\quad(1\leq k\leq n)$, \\

当$\displaystyle k=n$时,$\displaystyle x_{n,n}=\sqrt n\leq\sqrt n+1$.若$\displaystyle x_{k+1,n}\leq\sqrt{k+1}+1$, \\

则由$\displaystyle\sqrt{k+1}\leq2\sqrt k$可得$\displaystyle x_{k,n}\leq\sqrt{k+\sqrt{k+1}+1}\leq\sqrt{k+2\sqrt k+1}=\sqrt k+1,$ \\

故由反向归纳法,$\displaystyle a_n=x_{1,n}\leq2$,即$\displaystyle\{a_n\}$有上界.\\

又$\displaystyle\{a_n\}$单调递增,所以由单调收敛定理,$\displaystyle\{a_n\}$收敛. \\

\begin{ti}
【题5】设已知$\displaystyle\lim_{n\to\infty}a_n=a\quad$,证明$\displaystyle\lim_{n\to\infty}\frac{1}{2^n}\sum_{k=0}^{n}\binom{n}{k}a_k=a $
\end{ti}

\textbf{证明}:$\;\because 2^n=\displaystyle\sum_{k=0}^{n}\binom{n}{k},$ \\

$\therefore\left\lvert \dfrac{1}{2^n}\displaystyle\sum_{k=0}^{n}\binom{n}{k}a_k-a \right\rvert =\left\lvert \dfrac{1}{2^n}\displaystyle\sum_{k=0}^{n}\binom{n}{k}(a_k-a) \right\rvert$ $\leq\dfrac{1}{2^n}\displaystyle\sum_{k=0}^{n}\binom{n}{k}\left\lvert a_k-a\right\rvert,$ \\

$\because\displaystyle\lim_{n\to\infty}a_n=a\;\therefore\text{对}\forall\varepsilon>0,\;\exists N_1\;\text{s.t.}\;n>N_1\text{时},\;\left\lvert a_n-a\right\rvert<\varepsilon,$ \\

$\therefore\dfrac{1}{2^n}\displaystyle\sum_{k=0}^{n}\binom{n}{k}\left\lvert a_k-a\right\rvert=\dfrac{1}{2^n}\displaystyle\sum_{k=0}^{N_1}\binom{n}{k}\left\lvert a_k-a\right\rvert+\dfrac{1}{2^n}\displaystyle\sum_{k=N_1+1}^{n}\binom{n}{k}\left\lvert a_k-a\right\rvert$ \\

$<\dfrac{1}{2^n}\displaystyle\sum_{k=0}^{N_1}\binom{n}{k}\left\lvert a_k-a\right\rvert+\dfrac{\varepsilon}{2^n}\displaystyle\sum_{k=N_1+1}^{n}\binom{n}{k}$ \\

$<\dfrac{1}{2^n}\displaystyle\sum_{k=0}^{N_1}\binom{n}{k}\left\lvert a_k-a\right\rvert+\varepsilon,$ \\

又$\displaystyle\binom{n}{k} =\dfrac{n!}{k!(n-k)!}\leq\dfrac{n!}{(n-k)!}\leq n^k,\;\text{令}M=\max\{\left\lvert a_0-a\right\rvert,\left\lvert a_1-a\right\rvert,...,\left\lvert a_{N_1}-a\right\rvert  \},$ \\

$\therefore\dfrac{1}{2^n}\displaystyle\sum_{k=0}^{N_1}\binom{n}{k}\left\lvert a_k-a\right\rvert\leq\dfrac{M}{2^n}\displaystyle\sum_{k=0}^{N_1}n^k<\dfrac{M\cdot n^{N_1+1}}{2^n},$ \\

$\because\displaystyle\lim_{n\to\infty}\dfrac{n^{\alpha}}{2^n}=0\;\therefore\forall\varepsilon,\;\exists N_2,\;\text{s.t.}\;n>N_2\text{时},\dfrac{n^{N_1+1}}{2^n}<\varepsilon,$ \\

$\therefore\forall\varepsilon>0,\;\exists N=\max\{N_1,N_2\},\; n>N\text{时},$ $\left\lvert \dfrac{1}{2^n}\displaystyle\sum_{k=0}^{n}\binom{n}{k}a_k-a \right\rvert<(M+1)\varepsilon,$ \\

$\therefore\displaystyle\lim_{n\to\infty}\frac{1}{2^n}\sum_{k=0}^{n}\binom{n}{k}a_k=a.$

\begin{ti}
【题6】设${a_n}$满足条件$0<a_1<1,a_{n+1}=a_n(1-a_n)$,证明:$\displaystyle\lim_{n\to\infty}na_n=1$
\end{ti}

\textbf{证明}:$\;\because a_{n+1}-a_n=-a_n^2\leq 0,\quad\therefore \{a_n\}$单调递减, \\

由数学归纳法可得$0<a_n<1,\quad\therefore \{a_n\}$收敛, \\

设$\displaystyle\lim_{n\to\infty}a_n=\alpha,\text{其中}0\leq\alpha\leq 1,$ $\quad\therefore\alpha=\alpha(1-\alpha)\Rightarrow \alpha=0,$ \\

又$\dfrac{1}{a_{n+1}}=\dfrac{1}{a_n(1-a_n)}=\dfrac{1}{a_n}+\dfrac{1}{1-a_n},\quad\therefore\dfrac{1}{a_{n+1}}-\dfrac{1}{a_n}=\dfrac{1}{1-a_n},$ \\

$\therefore\displaystyle\lim_{n\to\infty}\dfrac{(n+1)-n}{\dfrac{1}{a_{n+1}}-\dfrac{1}{a_n}}=\displaystyle\lim_{n\to\infty}(1-a_n)=1,$ \\

又$\dfrac{1}{a_n}$为严格单调递增的正无穷大量, \\

$\therefore\displaystyle\lim_{n\to\infty}\dfrac{n}{\dfrac{1}{a_n}}=1,\;\text{即}\displaystyle\lim_{n\to\infty}na_n=1.$ \\

\textbf{注}:本题与题2类似,直接用Stolz公式并不好做,取倒数的动机为如果$\displaystyle\lim_{n\to\infty}na_n=1$ 

这就说明$a_n$的衰减速度与$\dfrac{1}{n}$相同,则$b_n=\dfrac{1}{a_n}$就是线性增长,便于研究.

\begin{ti}
【题7】证明:自然对数的底$e$是无理数
\end{ti}

\textbf{证明}:$\;\text{假设}e\in\mathbb{Q},\text{则}\exists\;p,q\in\mathbb{N}^+,\;\text{s.t.}\,e=\dfrac{p}{q},$ \\

设$\varepsilon_n=e-(1+\dfrac{1}{1!}+\dfrac{1}{2!}+...+\dfrac{1}{n!}),$ \\

先证明一个引理:$\;\varepsilon_n<\dfrac{1}{n!n},$ \\

$n!\varepsilon_n=n!\displaystyle\sum_{k=n+1}^{\infty}\dfrac{1}{k!}=\dfrac{1}{n+1}+\dfrac{1}{(n+1)(n+2)}+...,$ \\

$\dfrac{1}{n}-\dfrac{1}{n+1}=\dfrac{1}{n(n+1)},\dfrac{1}{n(n+1)}-\dfrac{1}{(n+1)(n+2)}=\dfrac{2}{n(n+1)(n+2)}...,$  \\

这样归纳下去,可得$\forall m\in\mathbb{N}^+ ,\dfrac{1}{n}>\dfrac{1}{n+1}+\dfrac{1}{(n+1)(n+2)}+...+\dfrac{1}{(n+1)(n+2)...(n+m)}\rlap{,}$ \\

从而$\varepsilon_n<\dfrac{1}{n!n},$\\

$\because q\in\mathbb{N}^+,\,e=(1+\dfrac{1}{1!}+...+\dfrac{1}{q!})+\varepsilon_q=\dfrac{p}{q},$ \\

两边同乘$\displaystyle q!$,有$\displaystyle p(q-1)!-q!\sum_{k=0}^{q}\dfrac{1}{k!}=q!\varepsilon_q,$ \\

等式左边为整数,故$\displaystyle q!\varepsilon_q\in\mathbb{Z}$.另一方面,$\displaystyle 0<q!\varepsilon_q<\dfrac{1}{q}\leq1$,即$\displaystyle 0<q!\varepsilon_q<1$,矛盾!

这就说明$e$是无理数,得证.

\begin{ti}
【题8】设已知存在极限$\displaystyle\lim_{n\to\infty}\dfrac{a_1+a_2+...+a_n}{n},$证明:$\displaystyle\lim_{n \to \infty}\dfrac{a_n}{n}=0  $
\end{ti}

\textbf{证明}: 设$b_n=\dfrac{a_1+a_2+...+a_n}{n},\text{由题意,不妨设}\{b_n\}\text{的极限为}l,$ \\

$\because a_n=nb_n-(n-1)b_{n-1}\,(n\geq 2),$ \\

$\therefore\displaystyle\lim_{n \to \infty}\dfrac{a_n}{n}=\lim_{n\to\infty}\dfrac{nb_n-(n-1)b_{n-1}}{n}=l-l=0.$

\begin{ti}
【题9】 设$x_n=\dfrac{1}{n^2}\displaystyle\sum_{k=0}^{n}ln\binom{n}{k},\text{求}\lim_{n\to\infty}x_n$
\end{ti}

\textbf{证明}: 由于$n^2$是单调递增的正无穷大量, \\

由Stolz公式,$\displaystyle\lim_{n\to\infty}x_n=\lim_{n\to\infty}\dfrac{\displaystyle\sum_{k=0}^{n+1}ln\binom{n+1}{k}-\displaystyle\sum_{k=0}^{n}ln\binom{n}{k}}{2n+1}$ \\

$=\displaystyle\lim_{n\to\infty}\dfrac{\displaystyle\sum_{k = 0}^{n}ln\frac{n+1}{n+1-k}}{2n+1},$又由于$2n+1$也是单调递增的正无穷大量, \\

再由Stolz公式,$\displaystyle\lim_{n\to\infty}\dfrac{\displaystyle\sum_{k = 0}^{n}\ln\frac{n+1}{n+1-k}}{2n+1}=\lim_{n\to\infty}\dfrac{\displaystyle\sum_{k=0}^{n}\ln\frac{n+1}{n+1-k}-\displaystyle\sum_{k=0}^{n-1}\ln\frac{n}{n-k}}{2}$ \\

$=\displaystyle\lim_{n\to\infty}\dfrac{\displaystyle\sum_{k=1}^{n}\ln\frac{n+1}{n+1-k}-\displaystyle\sum_{k=0}^{n-1}\ln\frac{n}{n-k}}{2}$ \\

$=\displaystyle\lim_{n\to\infty}\dfrac{\displaystyle\sum_{k=0}^{n-1}\ln\frac{n+1}{n-k}-\displaystyle\sum_{k=0}^{n-1}\ln\frac{n}{n-k}}{2}$ \\

$=\displaystyle\lim_{n\to\infty}\dfrac{n\ln\dfrac{n+1}{n}}{2},$ \\

$\because\dfrac{1}{n+1}\leq\ln(1+\dfrac{1}{n})\leq\dfrac{1}{n},\quad\therefore\dfrac{n}{n+1}\leq n\ln(1+\dfrac{1}{n})\leq 1,$ \\

由夹逼准则,$\displaystyle\lim_{n\to\infty}n\ln(1+\dfrac{1}{n})=1,$ \\

$\therefore\displaystyle\lim_{n\to\infty}x_n=\dfrac{1}{2}.$

\begin{ti}
【题10】设$a_n>0,\{\dfrac{a_{n+1}+a_{n+2}}{a_n}\}\text{收敛于}\alpha,\alpha>2$,证明:数列$\{a_n\}$没有上界
\end{ti}

\textbf{证明}:用反证法,假设$\{a_n\}$有上界,又$a_n>0\Rightarrow\{a_n\}$为有界序列, \\

所以$\{a_n\}$上下极限存在,记$a=\displaystyle\varlimsup_{n\to\infty}a_n,b=\displaystyle\varliminf_{n\to\infty}a_n,\text{那么}a\geq b\geq 0,$ \\

\ding{172}$\;a>0$:$\;\exists\,\beta>2,N_1,\;\text{s.t.}\,n>N_1\text{时},a_{n+1}+a_{n+2}>\beta a_n,$ \\

由上极限的保不等式性质,$2a\geq\beta a\Rightarrow\beta\leq 2,$矛盾!

\ding{173}$\;a=0$:同上有$a_{n+1}+a_{n+2}>\beta a_n(\beta>2,n>N_1),$ \\

记$M_n=\displaystyle\sup_{k\geq n}a_k,\text{则}a_{n+1}\leq M_n,a_{n+2}\leq M_n\Rightarrow M_n>\dfrac{\beta}{2}a_n>a_n,$ \\

又$M_n=\max\{a_n,M_{n+1}\},\;\therefore\,M_n=M_{n+1},\forall\,n>N_1,$ \\

又$a=0\Rightarrow\displaystyle\lim_{n\to\infty}M_n=0\;\therefore\forall\,n>N_1,\text{都有}M_n=0$,这与$a_n>0$矛盾! \\

综上,$\{a_n\}$没有上界. \\

\textbf{注意}:上下极限没有乘除法的四则运算性质,所以即使是在$a>b>0$的情况下,\\
$\alpha\neq\dfrac{2a}{b},\alpha\neq\dfrac{2b}{a}$ 



\BookEnd
